% !PRINTED 507
\weberpage{507}
6. \gesperrt{If $\omega$ is a functional in $\Omega$, and if there are $m$ quantities $\omega_1, \omega_2, \ldots, \omega_m$, which do not all vanish, of such a kind that the $m$ products $\omega\omega_i$ for $i=1,2,\ldots,m$ can be put into the form:}
\weq{5}{\omega\,\omega_i = A_{1,i}\omega_1 + A_{2,i}\omega_2 + \cdots + A_{m,i}\omega_m,}
\gesperrt{in which the $m^2$ symbols $A_{k,i}$ are integral rational functionals, then $\omega$ is an integral functional.}

Here $m$ is any natural integer. In the application the $\omega_i$ are always numbers or functionals in $\Omega$; for the validity of the theorem, however, only the practicability of the multiplication indicated in (5) is essential.

The theorem is a simple consequence of the definition of the integral functionals; for since the $\omega_i$ do not all vanish, by the determinant theorem (Bd. I, §.\,24, II.) we must have
\weq{6}{\begin{vmatrix}
A_{1,1}-\omega, & A_{2,1}, & \ldots, & A_{m,1} \\
A_{1,2}, & A_{2,2}-\omega, & \ldots, & A_{m,2} \\
\hdotsfor{4} \\
A_{1,m}, & A_{2,m}, & \ldots, & A_{m,m}-\omega
\end{vmatrix} = 0}
By arranging according to powers of $\omega$ there results from this an equation
\weq{7}{\omega^m + A_1\omega^{m-1} + \cdots + A_m = 0,}
in which the coefficients $A_1, A_2, \ldots$ are composed from the $A_{i,k}$ by addition and multiplication, and are therefore, by 1., integral rational functionals; accordingly $\omega$ is also an integral functional.

For the last coefficient $A_m$ in the equation (7) we obtain the expression
\weq{8}{(-1)^mA_m = \Sigma \pm A_{1,1}A_{2,2}\ldots A_{m,m},}
\gesperrt{which is therefore equal to the product of all the roots of the equation (7).}

Theorem 4. is important as a criterion for integral functionals; it serves us here for the proof of the following theorem: % [sic: Weber prints Theorem 4.; the determinant criterion just proved is Theorem 6.]

7. \gesperrt{If $\Psi(x,y,\ldots)$ is an integral function whose coefficients are integral rational numbers or functionals, but free of the variables $x,y,\ldots$, and if further $\alpha,\beta,\ldots$ are integral functionals in $\Omega$, then}
\weq{9}{\omega = \Psi(\alpha,\beta,\ldots)}
\gesperrt{is also an integral functional.}
